% Online Appendix to Section 10 (sec:network). Numbers: paper/results/network.json (final rerun of
% 2026-10-12, scripts/30 at lambda = 0.78). Specification: R/indices/allied_network_design.R,
% scripts/30_build_allied_network_design.R.
\section{The allied network design: further results}\label{oa:network}

\paragraph{Inputs.} The design reads five collapsed inputs restricted to the 63 economies of the
network: 2{,}079 $\EO$ and 1{,}881 $\ES$ records, 726 development-potential records, and 128{,}898
friendshore and 116{,}622 opportunity dyads. For each, the number of rows read equals the number of
keys kept, with no duplicated or multivalued key, so the collapse rules (one record per member,
technology and stage; one record per ordered pair and stage; the latest year and the mean of the
distinct finite values otherwise) act as guards only. All 37 nodes have an $\EO$ record, and hence a
producer score, in every one of the 22 stages. In Geothermal Upstream and Midstream all 37 scores
rest on $\EO$ alone, and in Green Hydrogen Midstream the scores of Costa Rica, the Dominican
Republic, Ghana, Guatemala and Namibia do. Without $\ES$, demand rests on size alone, so Namibia,
the smallest node, has zero demand in these three stages; it remains a potential producer.

\paragraph{Producers.} Eight nodes produce only Upstream, all of them resource economies
(Australia, Chile, Saudi Arabia, Peru, Argentina, Oman, Namibia and the Democratic Republic of the
Congo), while Japan takes ten of its twelve roles and Korea nine of its ten in Midstream. Canada
(16 roles, nine of them Upstream), the EU node (15), the United States and Indonesia (13 each) and
India and Japan (12 each) are followed by Thailand and Mexico (11 each). The largest share of a stage is at most
0.30 in 18 of the 22 stages, and 30\% of all shares sit at the floor of 0.05. The partner tier
carries 58\% of designed Upstream and 45\% of Midstream production. Indonesia leads Batteries, Coal,
Electric Vehicles, Green Hydrogen, Solar and Wind Upstream with shares of 0.225 to 0.325, and Oil
Midstream.

\paragraph{The EU node.} The EU node produces in 15 of the 22 stages (six Upstream, nine
Midstream), with shares from 0.05 to 0.215 and a median of 0.105. It is the second-largest producer
in Green Hydrogen Midstream (0.215) and Nuclear Midstream (0.205), and its 0.205 in Coal Midstream
equals Japan's leading share, but it is the sole largest producer in no stage. It is a top-5
producer in ten stages, exactly its cap, and the portfolio loop bars it from Electric Grid Upstream,
Electric Vehicles Midstream, Nuclear Upstream, Solar Midstream and Wind Upstream, so Electric
Vehicles and Solar are the two Midstream stages without EU production. By its extra-EU exports the
EU is the largest exporting node in nine of the 18 basket stages, all Midstream. The design sees it
differently in two ways: its producer score averages strong and weak members, and its demand is
rank-compressed, with a median weight of 0.035 across stages (0.029 to 0.046) against 24.6\% of the
nodes' GDP, so that it consumes about as much as the United States (35.6\% of the nodes' GDP, median
weight 0.033). Table~\ref{tab:network_eu_split} splits the EU's designed shares back to its members
in proportion to the stage weights that aggregate their inputs; summed over stages, Germany (0.41),
France (0.24), Poland (0.19) and Italy (0.16) carry the largest derived shares. These figures are
derived, not designed: the program never sees the members separately.

% Generated by scripts/paper/fig_network.R from data/processed/outputs/allied_network_member_shares.csv
\begin{table}[htbp]
\centering
\begin{threeparttable}
\caption{The EU node's designed production shares and their split to member states}
\label{tab:network_eu_split}
\footnotesize
\begin{tabular}{@{}l l r l l r@{}}
\toprule
Technology & Stage & EU share & Weights & Three largest members (derived share) & Top-3 weight \\
\midrule
Green Hydrogen & Upstream & .125 & GDP & DEU~.030, FRA~.020, ITA~.015 & .52 \\
Solar & Upstream & .115 & exports & FRA~.053, DEU~.020, NLD~.012 & .74 \\
Geothermal & Upstream & .085 & GDP & DEU~.020, FRA~.013, ITA~.010 & .52 \\
Batteries & Upstream & .055 & exports & DEU~.016, FRA~.006, FIN~.005 & .49 \\
Coal & Upstream & .050 & production & POL~.025, DEU~.014, CZE~.006 & .91 \\
Electric Vehicles & Upstream & .050 & exports & FIN~.014, DEU~.007, NLD~.007 & .54 \\
Green Hydrogen & Midstream & .215 & GDP & DEU~.051, FRA~.034, ITA~.026 & .52 \\
Nuclear & Midstream & .205 & exports & FRA~.058, DEU~.054, NLD~.028 & .68 \\
Coal & Midstream & .205 & exports & POL~.082, DEU~.032, ITA~.022 & .66 \\
Wind & Midstream & .155 & exports & DEU~.045, ITA~.016, POL~.012 & .47 \\
Gas & Midstream & .135 & exports & DEU~.038, ITA~.025, FRA~.010 & .54 \\
Batteries & Midstream & .105 & exports & POL~.023, HUN~.022, DEU~.020 & .62 \\
Electric Grid & Midstream & .095 & exports & DEU~.031, ITA~.008, FRA~.006 & .47 \\
Geothermal & Midstream & .070 & GDP & DEU~.017, FRA~.011, ITA~.008 & .52 \\
Oil & Midstream & .055 & exports & DEU~.016, ITA~.008, FRA~.004 & .51 \\
\bottomrule
\end{tabular}
\begin{tablenotes}[flushleft]\footnotesize
\item \emph{Notes:} Stages in which the EU-27 node produces in the production design (15 of 22). EU share: the node's designed production share. The design solves the EU as one node; its share is split back to the 27 member states in proportion to the stage weights that aggregate their inputs (Weights: the members' production or export shares in the stage, else their GDP shares), so the member figures are derived, not designed. Top-3 weight: the combined weight of the three largest members within the node. ISO3 codes. Source: \texttt{allied\_network\_member\_shares.csv}.
\end{tablenotes}
\end{threeparttable}
\end{table}


\paragraph{Flows and edges.} In Batteries Midstream the United States carries 0.295 of production,
India and Malaysia 0.135 each, Canada 0.114 and the EU 0.105, with four more nodes near the floor;
in Solar Midstream India (0.275), the United States (0.205), Viet Nam (0.135) and Thailand (0.125)
lead; in Wind Midstream India (0.275), the United States (0.250) and the EU (0.155). Of the 811
consumer--stage pairs with positive demand, 793 draw on exactly two suppliers and 18 on three.
Self-supply is positive in two stages only, 0.8\% of demand in Electric Grid Upstream and 0.2\% in
Electric Vehicles Midstream, because a node's own edge is worth the median edge while the edges to
its best partners lie above it. The flow-weighted mean edge weight over flows between distinct
nodes ranges from 0.385 to 0.910 across stages. The low end is Geothermal, whose edges lack the
opportunity component and so enter at no more than half weight (flow-weighted 0.385 to 0.387),
against 0.73 to 0.75 in the two Green Hydrogen stages and 0.85 to 0.91 elsewhere; the flow-weighted
mean exceeds the unweighted mean in all 22 stages. The EU node's edges are weighted by its members'
2024 basket trade with the partner where it exists (76\% of the EU edges) and by the stage weights
otherwise.

\paragraph{Concentration and current exporters.} The design is less concentrated than the nodes'
exports in 10 of the 18 basket stages and more in eight; the spread of concentration is compressed,
from 0.095--0.355 in the nodes' exports to 0.126--0.209 in the design. Where the nodes' exports are
most concentrated, in Nuclear Upstream (Canada 54\%), Coal Upstream (Australia 41\%) and Nuclear
Midstream (the EU 48\%), the design spreads production; where they are dispersed, as in Batteries
Midstream ($\HHI^{\mathrm{N}}=0.095$), it concentrates it on the highest-scoring producers. The
largest designed share is below the largest node's export share in 13 of the 18 stages, and the
designed and observed leaders coincide in four. Against world exports the design is more
concentrated in 13 of the 18 stages, but so are the nodes' own exports in 16. China leads eight of
the nine Midstream baskets (Nuclear Midstream, led by France, is the exception) and the Solar and
Wind Upstream baskets (silicon metal, 38\%, and magnet rare earths, 30\%); Kazakhstan leads Nuclear
Upstream (44\%). The number of exporters behind each index differs (34 to 139 world reporters, at
most 37 nodes, seven to ten designed producers), so part of any difference reflects the number of
units. All three largest exporting nodes are designed producers in eight of the 18 stages and at
least one is in every stage. The lowest correlations between designed and observed shares, in
Batteries and Electric Vehicles Upstream (0.27 and 0.29), illustrate the combinatorial logic: in
Batteries Upstream Indonesia (0.235) is placed level with Australia, the largest exporter, and
Canada and Chile receive 0.12 each, while South Africa and the Philippines, the second and third
exporters, receive 0.05 and 0.06.

\paragraph{Endowment.} Applying the endowment screen of the allied picks after the fact
(technology-level reserves or production above zero; the EU node counts as endowed when any member
is), the 21 roles held by nodes without a recorded resource base carry 5\% of Batteries Upstream
(the United Kingdom), 10\% of Nuclear Upstream (Israel and Thailand at the floor), 12.5\% of Coal
Upstream (Peru and the Philippines), 23.5\% of Wind Upstream (Japan, Israel, Korea and Singapore),
27.5\% of Electric Grid Upstream (five nodes at or near the floor), 28\% of Green Hydrogen Upstream
(Thailand 0.145, Israel and the United Kingdom) and 36\% of Solar Upstream (the United Kingdom,
Norway, South Africa and Singapore), and none of Electric Vehicles, Gas and Oil Upstream. These
nodes score on the stage through its trade basket, readiness and the other categories of $\EO$ and
$\ES$.

\paragraph{Portfolio caps.} Eleven nodes have no top-5 appearance. The final pass bars Thailand and
the United States from nine stages each and Indonesia and Malaysia from eight. Without caps, six
nodes would exceed their cap: the United States (17 appearances against a cap of ten), Indonesia
(12 against seven), India (ten against eight), Thailand and Malaysia (ten against six) and Canada
(eight against seven). The caps pass these appearances chiefly to Japan and Mexico (from three to
six each), the Philippines (none to four), South Africa (one to four) and Israel (two to four); 17
stages lose objective value. Four Upstream stages (Coal, Green Hydrogen, Nuclear and Solar) have a
tie at fifth place, which the strict rank resolves by producer score.

\paragraph{Build candidates.} The build priority is
\begin{equation}\label{eq:net_build}
  b_{p,ts}=\delta_{p,ts}\,\big(0.6\,a_{p,ts}+0.4\,\bar W_{p,ts}\big),\qquad
  \bar W_{p,ts}=\textstyle\sum_cW_{pc}\,d_c,
\end{equation}
development potential times a weighted mean of the producer score and the demand-weighted
connectivity. Restricted to the 22 nodes with an observed development potential, the ten leading
candidates are India in Solar Midstream (0.76), Wind Midstream (0.75), Gas Midstream (0.72) and
Nuclear Midstream (0.69); Indonesia in Electric Vehicles Upstream and Coal Upstream (0.74 each),
Batteries Upstream (0.71), Solar Upstream (0.69) and Wind Upstream (0.68); and Viet Nam in Solar
Midstream (0.67). All ten are already designed producers, as the construction implies: the priority
weights the quantities on which the program itself selects producers, so the list shows where the
design already relies on an emerging economy rather than where new capacity would add a producer.
